\section{An evaluated global affine pencil and its proved scope}
\label{sec:c160-global}

This section supplies actual rational matrices in the original eight
coordinates.  It proves that their projection contains the quantum set.
It does not prove the reverse inclusion.  Consequently, the dimensions
below are dimensions of an evaluated \emph{outer candidate}; they are not
an upper bound for an exact lift of $\mathcal Q_{222}$.

\paragraph{The canonical physical domain.}
Write $z=(k,a_0,b_0,a_1,b_1)$ and let $T\subset[0,1]^5$ be defined by
\[
\begin{split}
 g_1&=kb_0(1-a_0)-a_0(1-b_0)\ge0,\\
 g_2&=ka_1(1-b_0)-b_0(1-a_1)\ge0,\\
 g_3&=kb_1(1-a_1)-a_1(1-b_1)\ge0.
\end{split}\tag{CC1}
\]
Here the letters $a_x,b_y$ in $z$ are angular parameters; the visible
expectation coordinates will be denoted by $q$ to avoid confusing them
with these parameters.  Put
\[
 d(t)=(1-t)^2+t^2,\quad c(t)=\frac{1-2t}{d(t)},\quad
 s(t)=\frac{2t(1-t)}{d(t)},\quad
 D(z)=(1+k^2)\prod_{t\in\{a_0,b_0,a_1,b_1\}}d(t).
\tag{CC2}
\]
Since $c(t)^2+s(t)^2=1$, these are real qubit measurements
$c(t)Z+s(t)X$ on the normalized state
$(|00\rangle+k|11\rangle)/\sqrt{1+k^2}$.
Their eight-coordinate expectation map is
\[
\begin{split}
 \Phi_{A_x}(z)&=\frac{1-k^2}{1+k^2}c(a_x),\qquad
 \Phi_{B_y}(z)=\frac{1-k^2}{1+k^2}c(b_y),\\
 \Phi_{E_{xy}}(z)&=c(a_x)c(b_y)+\frac{2k}{1+k^2}s(a_x)s(b_y).
\end{split}\tag{CC3}
\]
All endpoints of the cube are physical and
$1/16\le D\le2$.  In particular, no limiting division by a vanishing
physical denominator is used.

Let $\mathcal G$ be the 128 input, outcome, and party relabelings in the
CHSH scenario, and let $\ell_1,\ldots,\ell_{16}$ be the deterministic
local vectors.  The extremal characterization of Barizien and Bancal
\cite{barizien_bancal} gives
\[
 \mathcal Q_{222}=\operatorname{conv}
 \left(\{\ell_1,\ldots,\ell_{16}\}\cup
       \bigcup_{g\in\mathcal G}g\Phi(T)\right).
\tag{CC4}
\]
For clarity, the conversion of its full-alternance condition into CC1
is algebraic.  After a Schmidt interchange one has $0<k\le1$.
Write $h(t)=t/(1-t)$, so that $2\arctan h(t)$ is the measurement
angle.  In the alternating half-circle ordering, the largest and
smallest steered Alice half-angle tangents are respectively
$h(a_x)/k$ and $kh(a_x)$.  Full alternance is therefore equivalent to
\[
 h(a_0)\le kh(b_0),\qquad h(b_0)\le kh(a_1),\qquad
 h(a_1)\le kh(b_1),
\]
which gives CC1 on clearing the positive denominators.  Outcome
relabelings select the half-circle representatives; reversal or an
exchange of the alternating ordering is included in $\mathcal G$.
Taking the closed domain includes the limiting angle and Schmidt
endpoints.  Local extreme points are the explicitly retained
deterministic vectors.  Thus CC4 uses a physical parametrization and
the published extremal theorem, rather than an unproved completeness
claim for a Gram cone.

\paragraph{Literal fixed features.}
The coefficient space has dimension $3^5=243$, with tensor basis
\[
 p_e(z)=\prod_{i=0}^4 p_{e_i}(z_i),\qquad
 (p_0(t),p_1(t),p_2(t))=((1-t)^2,t(1-t),t^2),
 \quad e\in\{0,1,2\}^5.
\tag{CC5}
\]
The recovered compiler has 160 blocks.  For each block $b$, the
accompanying file \nolinkurl{C160_FIXED_FEATURES.json} prints a subset
$I_b\subset\{0,\ldots,4\}$ and an explicit polynomial weight $w_b$.
Every weight is a positive rational constant times a product of
$z_i$, $1-z_i$, and $g_1,g_2,g_3$; it is therefore nonnegative on $T$.
The corresponding fixed vector is
\[
 v_b(z)=\sqrt{\frac{w_b(z)}{D(z)}}\,
       \left(\prod_{i\in I_b}z_i^{\eta_i}(1-z_i)^{1-\eta_i}
         \right)_{\eta\in\{0,1\}^{I_b}}.
\tag{CC6}
\]
Its entries depend only on the test point.  The block orders sum to
499 and the largest order is 32.  All 2,140 upper-triangular feature
product identities were checked exactly over $\mathbb Q$ against the
recovered compiler, including every off-diagonal factor of two.

Let $A_{r,c}$ be the exact coefficient matrix mapping upper-triangular
Gram coordinates to the basis CC5, and let $T_{r,j}$ be the nine target
columns.  Independent expansion of CC2--CC3 verifies
\[
 \sum_r T_{r,0}p_r=D,\qquad
 \sum_r T_{r,j}p_r=-D\Phi_j\quad(1\le j\le8).
\tag{CC7}
\]
Write $L_0=T_{\cdot,0}^{T}$ and $L_j=-T_{\cdot,j}^{T}$ for $j\ge1$.
For a moment vector $Y\in\mathbb R^{243}$, define each symmetric block
by
\[
 [M_b(Y)]_{ij}=\sum_r\frac{A_{r,c(b,i,j)}}{2-\delta_{ij}}Y_r,
 \qquad i\le j.
\tag{CC8}
\]
The denominator $2-\delta_{ij}$ is essential.  The stored primal
off-diagonal coordinate contributes twice a raw feature product;
the adjoint matrix entry contributes that raw product once.

\paragraph{Canonical affine elimination.}
The $9\times243$ matrix $L$ has rank nine.  Its exact pivot columns
and inverse are evaluated by the script.  Solving
\[
 LY=(1,q)^T
\tag{CC9}
\]
eliminates nine moments and leaves 234 auxiliaries.  Substitution into
CC8 gives the actual rational affine pencil in
\nolinkurl{C160_CANONICAL_OUTER_PENCIL.coo.json}.  It contains 3,749
nonzero upper-triangular coefficient entries.  Its block-diagonal
matrix order is 499.  Define its projection to be $R_{160}$.

For $z\in T$, choose $Y_e=p_e(z)/D(z)$.  Equations CC7 and CC8 give
$LY=(1,\Phi(z))^T$ and
$M_b(Y)=v_b(z)v_b(z)^T\succeq0$.  Convex combinations give
\[
 \operatorname{conv}\Phi(T)\subseteq R_{160}.
\tag{CC10}
\]
This inclusion is exact and includes every boundary point of $T$.

\paragraph{Compactness and the homogeneous zero face.}
Block zero has all five variables free and weight one.  Its products
span the complete 243-dimensional space.  The target $D$ has the
exact strictly positive block-zero Gram
\[
 Q_D=\begin{pmatrix}I_{16}&I_{16}\\I_{16}&2I_{16}\end{pmatrix},
 \qquad \lambda_{\min}(Q_D)=\frac{3-\sqrt5}{2}.
\tag{CC11}
\]
Consequently
$L_0Y=\operatorname{tr}(Q_DM_0(Y))\ge0$ whenever all the blocks
are positive semidefinite.  If $L_0Y=0$, positive definiteness of
$Q_D$ forces $M_0(Y)=0$, and the spanning property then forces
$Y=0$.  If $L_0Y=1$, the trace of $M_0(Y)$ is at most
$(3+\sqrt5)/2$.  Each $Y_e$ is one of its entries, so the feasible
moment set is bounded as well as closed.  Its linear image $R_{160}$
is compact.  This also proves that homogenization introduces no
nonzero moment vector of zero mass.

\paragraph{A single global affine pencil.}
Introduce a moment vector $Y_g\in\mathbb R^{243}$ for each of the
128 relabelings and local coefficients $\lambda_\ell\ge0$.  Impose
\[
\begin{split}
 M_b(Y_g)&\succeq0\quad(1\le b\le160,\ g\in\mathcal G),\\
 \sum_gL_0Y_g+\sum_\ell\lambda_\ell&=1,\\
 q&=\sum_g g(L_1Y_g,\ldots,L_8Y_g)^T
       +\sum_\ell\lambda_\ell\ell.
\end{split}\tag{CC12}
\]
The cached Bell matrices act on \emph{support directions}; their
transposes are used for $g$ on primal points.  All 128 matrices were
checked to be distinct signed permutations preserving the local
vertex set.  The lifted local vectors $(1,\ell)^T$ span
$\mathbb R^9$.  Eliminate nine $\lambda_\ell$ through an evaluated
rational $9\times9$ inverse and retain the other seven as auxiliaries.
The nine eliminated weights retain their scalar PSD constraints.

The resulting single block-diagonal pencil is fully printed in
\nolinkurl{C160_GLOBAL_OUTER_PENCIL.coo.json}, with
\[
 \boxed{N=128\cdot499+16=63\,888,\qquad
        m=128\cdot243+7=31\,111.}
\tag{CC13}
\]
It has 20,496 blocks, largest order 32, and 455,750 nonzero rational
upper-triangular coefficient entries.  The coordinate order is
$(a_0,a_1,b_0,b_1,E_{00},E_{01},E_{10},E_{11})$, now denoting the
original expectations.  Each stored row has the form
\[
 (\text{block},\text{row},\text{column},\text{variable},
   \text{numerator},\text{denominator}).
\]
Variable zero is the constant; variables 1--8 are the expectations;
the remaining variables are the 31,111 auxiliaries.  Each off-diagonal
entry is mirrored once.  Thus the file gives actual matrices
$A_0,A_1,\ldots,A_8,B_1,\ldots,B_m$, rather than a symbolic
instruction to search for them.

Let $R_{\mathrm{global}}$ be this projected spectrahedron.  CC11
shows directly that CC12 is the convex hull of the relabeled copies
of $R_{160}$ and the sixteen local points.  Together with CC4 and
CC10 it proves
\[
 \boxed{\mathcal Q_{222}\subseteq R_{\mathrm{global}}.}
\tag{CC14}
\]
Every homogeneous mass is nonnegative, their sum is one, and
zero-mass copies vanish.  The same estimates prove compactness of
the global feasible moment set.  Hence CC14 does not depend on
replacing a projection by its closure.

\paragraph{The remaining assertion is mathematical, not a compiler step.}
Equality requires
\[
 \boxed{R_{\mathrm{global}}\subseteq\mathcal Q_{222}.}
\tag{CC15}
\]
Equivalently, for every $h\in\mathbb R^8$, the maximum of $h\cdot q$
over CC12 must equal its quantum support value.  It would suffice to
prove that every affine functional valid on $\mathcal Q_{222}$ has
its canonical numerator in the fixed C160 Gram cone, with compatible
relabelings.  Neither the exact compiler identities, the physical
Dirac-moment witnesses, compactness, nor numerical agreement for
chosen Bell functionals proves this universal assertion.  CC15
remains unproved in this extraction package.
